\section{Related Work}
\label{sec:related}

\paragraph{Training dynamics.}
Arpit~et~al.~\cite{arpit2017closer} showed networks fit simple patterns before
memorising noise; Toneva~et~al.~\cite{toneva2019empirical} identified
\emph{forgettable} examples as disproportionately informative; and Dataset
Cartography~\cite{swayamdipta2020dataset} formalised an \emph{ambiguous}
stratum along confidence/variability axes.
EWL recovers an online analogue of this stratum in the same training run, with
no preliminary pass.

\paragraph{Curriculum, self-paced, and noise-robust learning.}
Curriculum~\cite{bengio2009curriculum} and self-paced
learning~\cite{kumar2010self} weight by loss \emph{magnitude}, conflating
frontier with noise.
MentorNet~\cite{jiang2018mentornet}, Co-teaching~\cite{han2018coteaching},
OHEM~\cite{shrivastava2016ohem}, and Focal Loss~\cite{lin2017focal} either
require a clean reference set or rely on magnitude heuristics.
EWL instead uses loss \emph{velocity} --- the relative rate of decline --- which
distinguishes actively declining samples from those stagnating at high loss,
with no explicit noise model.

\paragraph{Meta-reweighting and LoRA.}
Ren~et~al.~\cite{ren2018learning} and Meta-Weight-Net~\cite{shu2019meta} learn
weights via second-order gradients and a clean held-out set, at roughly
$2\times$ compute.
EWL uses only first-order in-run statistics.
LoRA~\cite{hu2022lora} constrains weight updates to a rank-$r$ subspace;
AdaLoRA~\cite{zhang2023adalora} reallocates \emph{rank} across layers using
adapter \emph{weights}.
EWL uses adapter \emph{gradient} norms for a complementary purpose: gating
per-sample reweighting by current adapter activity.
